\clearpage
\section{Node Embeddings with TransE}

A multi-relational graph is \(G=(\mathcal E,\mathcal S,\mathcal L)\), where
\(\mathcal E\) is the entity set, \(\mathcal L\) is the relation set, and
\(\mathcal S\) contains triples \((h,\ell,t)\). TransE learns entity and
relation embeddings in \(\R^k\) so that
\[
  \vect{h}+\vect{\ell}\approx\vect{t}
\]
for observed triples.

For a positive triple, define its corruption set by
\[
  \mathcal S'_{(h,\ell,t)}
  =
  \left\{(h',\ell,t):h'\in\mathcal E\right\}
  \cup
  \left\{(h,\ell,t'):t'\in\mathcal E\right\}.
\]
The TransE margin loss is
\[
  \mathcal L
  =
  \sum_{(h,\ell,t)\in\mathcal S}
  \sum_{(h',\ell,t')\in\mathcal S'_{(h,\ell,t)}}
  \left[
    \gamma+
    d(\vect{h}+\vect{\ell},\vect{t})
    -d(\vect{h'}+\vect{\ell},\vect{t'})
  \right]_+,
\]
where \(d\) is Euclidean distance, \([a]_+=\max(0,a)\), and entity embeddings
are constrained by \(\norm{\vect e}_2=1\).

\subsection{Simplified Loss}

Suppose the objective is replaced by
\[
  \mathcal L_{\mathrm{simple}}
  =
  \sum_{(h,\ell,t)\in\mathcal S}
  d(\vect h+\vect\ell,\vect t).
\]
Give a non-trivial graph with at least two nodes and one edge, together with
two-dimensional embeddings, that achieves zero loss but is completely
uninformative.

\begin{solution}
Let
\[
  \mathcal E=\{a,b\},\qquad
  \mathcal L=\{r\},\qquad
  \mathcal S=\{(a,r,b)\}.
\]
Choose the two-dimensional embeddings
\[
  \vect a=\vect b=
  \begin{bmatrix}1\\0\end{bmatrix},
  \qquad
  \vect r=
  \begin{bmatrix}0\\0\end{bmatrix}.
\]
The entity embeddings satisfy the unit-norm constraint.  Moreover,
\[
  d(\vect a+\vect r,\vect b)
  =
  \norm{
    \begin{bmatrix}1\\0\end{bmatrix}
    +
    \begin{bmatrix}0\\0\end{bmatrix}
    -
    \begin{bmatrix}1\\0\end{bmatrix}
  }_2
  =0,
\]
so \(\mathcal L_{\mathrm{simple}}=0\).

This representation is useless because the two distinct entities have exactly
the same embedding.  In fact, for every \(x,y\in\{a,b\}\),
\[
  \vect x+\vect r=\vect y,
\]
so the embeddings predict every possible \(r\)-triple and cannot determine
which edge is actually present.
\end{solution}

\subsection{Removing the Margin}

Now remove \(\gamma\) and optimize
\[
  \mathcal L_{\mathrm{no\ margin}}
  =
  \sum_{(h,\ell,t)\in\mathcal S}
  \sum_{(h',\ell,t')\in\mathcal S'_{(h,\ell,t)}}
  \left[
    d(\vect h+\vect\ell,\vect t)
    -d(\vect{h'}+\vect\ell,\vect{t'})
  \right]_+.
\]
Give a non-trivial graph and two-dimensional embeddings that achieve zero
loss but do not reveal whether a pair of entities is connected by a relation.

\begin{solution}
Use the same graph
\[
  \mathcal E=\{a,b\},\qquad
  \mathcal L=\{r\},\qquad
  \mathcal S=\{(a,r,b)\},
\]
but choose
\[
  \vect a=\vect b=
  \begin{bmatrix}1\\0\end{bmatrix},
  \qquad
  \vect r=
  \begin{bmatrix}0\\1\end{bmatrix}.
\]
The relation embedding is now nonzero.  The positive-triple distance is
\[
  d(\vect a+\vect r,\vect b)=1.
\]
Every corrupted head or tail is either \(a\) or \(b\), and both have the same
embedding.  Hence every corrupted triple also has distance
\[
  d(\vect{h'}+\vect r,\vect{t'})=1.
\]
Each summand in the no-margin loss is therefore
\[
  [1-1]_+=0,
\]
so \(\mathcal L_{\mathrm{no\ margin}}=0\).

Nevertheless, all entity pairs receive the same score.  The objective only
requires a positive triple to be no farther away than a corrupted triple; in
the absence of \(\gamma\), equality is sufficient.  The positive margin is
what rules out this all-scores-tied solution.
\end{solution}

\subsection{Removing Entity Normalization}

TransE normally renormalizes every entity embedding to unit length. Describe
how the algorithm could trivially reduce the loss if this normalization step
were omitted, and characterize the resulting embeddings.

\begin{solution}
Consider any collection of unnormalized embeddings
\(\{\bar{\vect e}\}\) and \(\{\bar{\vect\ell}\}\).  Scale every entity and
relation by a common factor \(c>0\):
\[
  \vect e=c\bar{\vect e},
  \qquad
  \vect\ell=c\bar{\vect\ell}.
\]
For one positive triple and one of its corruptions, let
\[
  d_+
  =d(\bar{\vect h}+\bar{\vect\ell},\bar{\vect t}),
  \qquad
  d_-
  =d(\bar{\vect{h'}}+\bar{\vect\ell},\bar{\vect{t'}}).
\]
Euclidean distance is homogeneous, so the corresponding loss term after
scaling is
\[
  [\gamma+c d_+-c d_-]_+
  =
  [\gamma-c(d_--d_+)]_+.
\]
Whenever \(d_->d_+\), choosing
\[
  c\geq \frac{\gamma}{d_--d_+}
\]
makes this term zero.  For a finite training set whose positive--negative
gaps are all strictly positive, one sufficiently large \(c\) makes every
hinge term zero.

Thus the optimizer can satisfy the fixed margin partly by increasing the
scale of all embeddings instead of learning a better normalized geometry.
The entity and relation norms can grow without bound, points become
arbitrarily far from the origin, and absolute distances become meaningless
and numerically unstable.  Unit normalization removes this scale degree of
freedom and forces the model to improve relative directions and positions
within a bounded space.
\end{solution}

\subsection{Limits of TransE Expressiveness}

Give a simple multi-relational graph for which no perfect embedding exists:
there is no assignment satisfying
\[
  \vect u+\vect\ell=\vect v
  \quad\text{for all }(u,\ell,v)\in\mathcal S
\]
while also satisfying
\[
  \vect u+\vect\ell\ne\vect v
  \quad\text{for all }(u,\ell,v)\notin\mathcal S.
\]
Use two-dimensional embeddings and explain what the example shows about the
types of relations TransE can and cannot model.

\begin{solution}
Let
\[
  \mathcal E=\{a,b,c\},\qquad
  \mathcal L=\{r,s\},
\]
and let the only positive triples be
\[
  \mathcal S
  =
  \{(a,r,b),\ (a,r,c),\ (b,s,a)\}.
\]
In particular, \((c,s,a)\notin\mathcal S\).

If a perfect TransE embedding existed, the first two positive triples would
require
\[
  \vect a+\vect r=\vect b,
  \qquad
  \vect a+\vect r=\vect c.
\]
Therefore
\[
  \vect b=\vect c.
\]
The third positive triple requires
\[
  \vect b+\vect s=\vect a.
\]
Substituting \(\vect c=\vect b\) then gives
\[
  \vect c+\vect s=\vect a.
\]
This predicts \((c,s,a)\) as a positive triple, contradicting
\((c,s,a)\notin\mathcal S\).  Hence no perfect embedding exists in
\(\R^2\); in fact, the contradiction holds in every dimension.

The failure comes from the fixed-translation form of TransE.  For a given
head embedding and relation embedding, \(\vect h+\vect r\) is a single point.
A one-to-many relation such as
\[
  (a,r,b),\qquad(a,r,c)
\]
therefore forces all of its tails to share one embedding.  If those tails
must play different roles elsewhere in the graph, TransE necessarily creates
false triples.  Thus TransE naturally represents one-to-one translations but
cannot faithfully express general one-to-many, many-to-one, or many-to-many
relations while keeping the involved entities distinguishable.
\end{solution}
